\section{Label placement, printing and practical limits}
Labels support \key{outside}, \key{left}, \key{inside}, \key{end}, \key{auto}
and \key{none}. \key{auto} moves a label outside
when it does not fit. \key{end} places it above the endpoint. Comparisons, lollipops,
outlined bars and intervals with a range point move internal labels outside.
Sparklines always use outside labels. \key{text color=auto} chooses black or white.
\key{label width} reserves an outside slot; use the same sufficiently wide slot
throughout a column. Long labels may require a wider column or a custom short text.

\begin{demo}{Left labels and inline baselines}
\gradbar[label=left,unit={\%},baseline=text]{92.7}
\qquad
\gradbar[label=left,baseline=bottom]{NA}
\end{demo}
\key{label=left} places the formatted label immediately before the track and
right-aligns it within \key{label width}. It is intended for compact table columns
where the value should read before the visual bar. \key{baseline=text} uses the
current label font metrics; \key{baseline=center} and \key{baseline=bottom} are
useful when a table row deliberately aligns graphics by their visual center or
lower edge. A positive \key{baseline shift} moves the picture upward.

\key{row padding} reserves extra height above and below a graphic; it does not
change the document's global row spacing. Use arraystretch or table row spacing
when appropriate. Marker shapes reserve space consistently even when the first
or last observation does not have a marker.

\key{print} applies gray fills, textures and white label backgrounds. Mono is a
group palette using those settings. Provide textual category names and real values
as well as visual distinctions. Outlined bars and legends are available through
\key{outline=true} or the outline preset; print clears outline mode.

All collision handling is local to a graphic, not to the whole table. Complex
axes, date parsing and large general-purpose plots are outside this package's
scope. Automatic alternative descriptions are available when PDF tagging is enabled;
see the tagging notes in the interface reference.
